\section{Related Work}

\paragraph{Diffusion and Flow-based Models.}
Diffusion models~\cite{diffusion,ddpm,edm,ncsn,scoresde} and flow matching~\cite{fm,rectified,stochastic} lay the foundation for a series of modern generative methods. 
These approaches can be formulated as learning a probabilistic trajectory, \ie an ODE/SDE (ordinary/stochastic differential equation) that maps between distributions.
A network is trained to model the underlying trajectory using a regression loss.
Samples are generated by solving the resulting ODE or SDE, typically using a numerical solver.

\paragraph{Fastforward Generative Models.}
%
Standard diffusion and flow-based models were originally designed without explicitly considering the acceleration of ODE/SDE solving.
An emerging category of methods, which we abstract as ``fastforward generative models'', explicitly incorporates ODE/SDE acceleration into their training objectives.

These models typically operate by making large jumps across time steps.
Consistency Models~\cite{cm,ict,ect,scm} formulate it as leaping from an intermediate time step directly to the end point of the trajectory.
Consistency Trajectory Models~\cite{ctm} aim to learn a trajectory between any two time steps, based on explicit integration (\ie ODE/SDE solving) during training.
Flow Map Matching~\cite{fmm} formulates the regression of the zeroth- and first-order derivatives of these flow fields.
Shortcut Models~\cite{shortcut} are built on the relationship between two time steps and their midpoint.
IMM~\cite{imm} leverages moment matching at different time steps.
MeanFlow~\cite{mf} formulates and parameterizes the average velocity across two arbitrary time steps.


Several improvements have been made to the MeanFlow formulation.
AlphaFlow \cite{alphaflow} decomposes the MeanFlow objective and adopts a schedule to interpolate from Flow Matching to MeanFlow.
Decoupled MeanFlow \cite{dmf} fine-tunes pre-trained Flow Matching models into MeanFlow by conditioning the final blocks of the networks on a second timestep.
CMT \cite{cmt} introduces mid-training using fixed explicit regression targets supplied by a pre-trained Flow Matching model before training the fastforward models. 
Our iMF is focused on the fundamental limitations of the MeanFlow objective, as well as the practical issue of CFG. These issues are orthogonal to other concurrent improvements.
